\DocumentMetadata{
  pdfversion=2.0,
  pdfstandard=ua-2,
  lang=en-US,
  tagging=on,
  tagging-setup={math/setup=mathml-SE,tabsorder=structure}
}
\documentclass[11pt,letterpaper]{article}
\usepackage[margin=0.68in,includefoot]{geometry}
\usepackage{newcomputermodern}
\usepackage{amsmath,amscd,mathtools}
\providecommand{\square}{\mdlgwhtsquare}
\usepackage[hidelinks]{hyperref}
\hypersetup{
  pdftitle={Analysis PhD Exam, May 2025},
  pdfauthor={University of Florida Department of Mathematics},
  pdfsubject={Graduate mathematics examination},
  pdfdisplaydoctitle=true,
  bookmarksopen=true
}
\setcounter{secnumdepth}{0}
\setlength{\parindent}{0pt}
\setlength{\parskip}{0.28em}
\setlength{\emergencystretch}{3em}
\setlength{\leftmargini}{2.15em}
\setlength{\leftmarginii}{2.65em}
\setlength{\leftmarginiii}{3em}
\renewcommand{\labelenumi}{\textbf{\theenumi.}}
\pagestyle{empty}
\raggedbottom
\allowdisplaybreaks
\newenvironment{examproblems}[1]{%
  \begin{enumerate}%
  \setcounter{enumi}{#1}%
  \setlength{\topsep}{0.35em}%
  \setlength{\partopsep}{0pt}%
  \setlength{\itemsep}{0.48em}%
  \setlength{\parsep}{0.18em}%
}{\end{enumerate}}
\newenvironment{examsubparts}{%
  \begin{enumerate}%
  \setlength{\topsep}{0.18em}%
  \setlength{\partopsep}{0pt}%
  \setlength{\itemsep}{0.16em}%
  \setlength{\parsep}{0.1em}%
}{\end{enumerate}}
\newenvironment{examinstructions}{%
  \par\begingroup\itshape
}{%
  \par\endgroup\vspace{0.18em}
}
\makeatletter
\renewcommand\section{\@startsection{section}{1}{0pt}%
  {-1.25ex plus -.25ex minus -.1ex}%
  {0.55ex plus .1ex}%
  {\normalfont\large\bfseries}}
\renewcommand{\@maketitle}{%
  \newpage\null\vskip -1.2em%
  {\footnotesize\bfseries
    \href{https://www.ufl.edu/}{University of Florida}, \href{https://math.ufl.edu/}{Department of Mathematics}\par}%
  \vskip 0.18em%
  \tagstructbegin{tag=Title}%
    {\LARGE\bfseries\href{https://gma.math.ufl.edu/exams/analysis-phd/}{Analysis PhD Exam}, May 2025\par}%
  \tagstructend%
  \vskip 0.35em%
  \tagmcbegin{artifact}%
    \hrule height 1pt%
  \tagmcend%
  \vskip 0.55em%
}
\makeatother
\title{Analysis PhD Exam, May 2025}
\author{}
\date{}
\begin{document}
\maketitle
\thispagestyle{empty}
\begin{examinstructions}
Write solutions in a neat and logical fashion, giving complete reasons for all steps. Do six of the eight problems.
\end{examinstructions}
\begin{examproblems}{0}
\item
For any two of the following, either give an example or explain why no example exists.
\begin{examsubparts}
\item[\textnormal{(i)}]
A Lebesgue measurable set \(E \subseteq [0,1]\) that is closed, has positive measure, but contains no non-trivial interval.
\item[\textnormal{(ii)}]
A Lebesgue measurable set \(E \subseteq [0,1]\) of positive measure such that \(E-E=\{x-y:x,y\in E\}\) contains no non-trivial interval.
\item[\textnormal{(iii)}]
A measure space \((X,\mathcal M)\) and a set \(E\subseteq X\times X\) such that \([E]_x,[E]^y\in\mathcal M\) for all \(x,y\in X\), but \(E\notin\mathcal M\otimes\mathcal M\).
\end{examsubparts}
\item
Suppose \(X,Y\) are metric spaces. Show, if \(f:X\to Y\) is continuous, then~\(f\) is Borel measurable.
\item
Let \(I=(-1,1)\) and let \(m^*\) denote Lebesgue outer measure. Show, if \(A\subseteq I\) and \(m^*(A)+m^*(I\setminus A)=2\), then~\(A\) is Lebesgue measurable.
\item
State the Radon–Nikodym Theorem and give an example that shows the~\(\sigma\)-finite hypothesis is needed.
\item
Suppose \(\mathcal X\) is a Banach space and \(\mathcal M\) and \(\mathcal N\) are (closed) subspaces of \(\mathcal X\). Show, if each \(x\in\mathcal X\) has a unique representation as \(x=m+n\) for some \(m\in\mathcal M\) and \(n\in\mathcal N\), then the mapping \(\mathcal X\to\mathcal M\) sending \(x=m+n\) to~\(m\) is linear and bounded.
\item
Prove, if \(\mathcal X\) is a Banach space and \(\mathcal X^*\) is separable, then \(\mathcal X\) is separable. Is \(\ell^1\) isomorphic, as a Banach space, to \((\ell^\infty)^*\)?
\item
Short answer. Do two.
\begin{examsubparts}
\item[\textnormal{(i)}]
Suppose \(h:\mathbb R\to\mathbb C\) is a measurable function. If \(hf\in L^1(\mathbb R)\) for every \(f\in L^4(\mathbb R)\), what can be said about~\(h\)? Explain your answer.
\item[\textnormal{(ii)}]
Does there exist a (bounded) linear functional \(L:\ell^\infty(\mathbb N)\to\mathbb C\) such that \(L(f)=\lim f(n)\) whenever \(f\in\ell^\infty(\mathbb N)\) and \((f(n))_n\) converges? Explain your answer.
\item[\textnormal{(iii)}]
Does there exist an inner product on \(\mathbb R^2\) that induces the norm \(\lVert x\rVert=|x_1|+|x_2|\) for \(x=\binom{x_1}{x_2}\in\mathbb R^2\)? Explain your answer.
\end{examsubparts}
\item
Do two providing short justifications for your answers. Here \(\widehat{\phantom{f}}\) denotes the Fourier transform of an \(L^1\) function and \(\mathcal Ff\) denotes the Fourier transform of an \(L^2(\mathbb R)\) function~\(f\).
\begin{examsubparts}
\item[\textnormal{(i)}]
Define \(f\in L^1(\mathbb R)\) by \(f(x)=e^{-x^4}\chi_{[1,\infty)}(x)\). Is \(\widehat f\in L^1(\mathbb R)\)?
\item[\textnormal{(ii)}]
Recall, if \(f\in L^1(\mathbb R)\), then \(\lVert\widehat f\rVert_\infty\leq\lVert f\rVert_1\). When does equality hold?
\item[\textnormal{(iii)}]
Does there exist \(0\neq f\in L^2(\mathbb R)\) such that \(\mathcal Ff=e^{i\pi/4}f\)?
\end{examsubparts}
\end{examproblems}
\end{document}
